% BT1896 -- residual-language clarification patch for photonic_holonet.tex
% Suggested insertion: Appendix B, after the opening architecture-completeness paragraphs.

\subsection*{BT1896 Clarification: finite closure versus physical-continuum closure}

The residual ledger should be read in two layers.  The finite Holonet architecture is closed as a
machine specification: the $W(3,3)$ carriers, routing atlas, Steinberg memory, Witting transaction
ABI, contextual-fuel accounting, and finite lattice/register data are exact or machine-witnessed.
This is the sense in which the architecture is complete.

The physical-continuum interpretation is a second layer and remains deliberately qualified.  In
particular:

\begin{itemize}
\item \textbf{R1.}  The finite $E_8/2E_8$ quadratic refinement and positive $E_8$ lift are fixed as
finite gauge-lattice data.  This does not by itself constitute a completed continuum
Einstein--Hilbert derivation.
\item \textbf{R2.}  The moonshine traces and finite matter-register bridge are supplied at the
substrate level.  The physical fermion-multiplet assignment and full CKM/PMNS derivation remain
separate phenomenological identifications unless explicitly proved elsewhere.
\item \textbf{R3.}  The spacetime/field-theory limit is theorem-governed by known convergence and
holographic-code mechanisms, but the constructive physical continuum lift and its measured
Einstein--Hilbert coefficient remain a downstream target.
\end{itemize}

Thus the safe public statement is:
\[
\text{finite machine-complete architecture, with remaining physical/continuum identifications classified.}
\]
The stronger statement that all three physical residuals are fully solved should be avoided unless
and until the continuum and phenomenological derivations are inserted as independent theorems.
